\documentclass[12pt,a4paper]{article}
\usepackage[utf8]{inputenc}
\usepackage[T1]{fontenc}
\usepackage[spanish]{babel}
\usepackage{geometry}
\geometry{top=2.5cm, bottom=2.5cm, left=3cm, right=2.5cm}
\usepackage{longtable}
\usepackage{booktabs}
\usepackage{parskip}
\usepackage{hyperref}
\usepackage{xcolor}
\usepackage{array}
\usepackage{fancyhdr}
\usepackage{titlesec}

\definecolor{navyblue}{RGB}{0, 51, 102}
\titleformat{\section}{\Large\bfseries\color{navyblue}}{}{0em}{}[\vspace{-0.3em}\rule{\textwidth}{0.5pt}]
\titleformat{\subsection}{\large\bfseries\color{navyblue}}{}{0em}{}

\hypersetup{colorlinks=true, linkcolor=navyblue, urlcolor=navyblue}
\newcolumntype{L}[1]{>{\raggedright\arraybackslash}p{#1}}

\pagestyle{fancy}
\fancyhf{}
\fancyhead[L]{\small\color{navyblue}Términos de Referencia}
\fancyhead[R]{\small\color{navyblue}RENZO-TDR-2026-001}
\fancyfoot[C]{\thepage}
\renewcommand{\headrulewidth}{0.4pt}

\begin{document}

\thispagestyle{empty}
\begin{center}
\vspace*{2cm}
{\LARGE\bfseries\color{navyblue} TÉRMINOS DE REFERENCIA}

\vspace{0.5cm}
{\Large Contratación del Servicio de Instalación y Puesta en Servicio de\\
Sistema de Puesta a Tierra, Electrobomba\\
y Grupo Electrógeno (A Todo Costo)}

\vspace{1.5cm}

\rule{\textwidth}{1pt}
\vspace{0.5cm}

\begin{tabular}{rl}
\textbf{Proyecto:} & Instalaciones Eléctricas Interiores para Vivienda Unifamiliar de 3 Niveles \\
\textbf{Propietario:} & Renzo Gabriel Mamani Galindo \\
\textbf{Ubicación:} & Jr. Lima S/N, Capachica, Puno \\
\textbf{Código:} & RENZO-TDR-2026-001 \\
\textbf{Fecha:} & \today \\
\end{tabular}

\vspace{1.5cm}
\rule{\textwidth}{1pt}
\end{center}

\newpage
\tableofcontents
\newpage

% ===================================================================
\section{Finalidad pública}
% ===================================================================

Garantizar la seguridad eléctrica, la continuidad operativa y el abastecimiento hídrico de la vivienda unifamiliar de tres niveles ubicada en Jr. Lima S/N, Capachica, Puno. Mediante la implementación de un sistema de puesta a tierra eficiente, un sistema de bombeo de agua y un grupo electrógeno de respaldo, se busca proteger los equipos sensibles contra descargas y fallas, además de asegurar el funcionamiento ininterrumpido de los servicios críticos ante eventuales cortes del suministro comercial.

% ===================================================================
\section{Descripción general del requerimiento}
% ===================================================================

Contratación del servicio integral a todo costo (suministro, instalación, pruebas y puesta en marcha) de tres (03) componentes electromecánicos para la vivienda unifamiliar:

\begin{enumerate}
  \item Sistema de pozo a tierra (resistencia no mayor a 5 ohmios).
  \item Sistema de bombeo mediante electrobomba de 2 HP.
  \item Grupo electrógeno de respaldo dimensionado según máxima demanda.
\end{enumerate}

% ===================================================================
\section{Condiciones de contratación}
% ===================================================================

\begin{itemize}
  \item \textbf{Modalidad de pago:} Suma Alzada (Art. 130 del Reglamento de la Ley N° 32069).
  \item \textbf{Sistema de entrega:} Llave en mano. Ejecución integral y a todo costo (materiales, mano de obra especializada, herramientas, pruebas y certificación).
  \item \textbf{Plazo de entrega:} Máximo de 10 días calendario a partir del día siguiente de la notificación de la orden de servicio.
  \item \textbf{Lugar de prestación:} Jr. Lima S/N, Capachica, Puno. Horario: Lunes a Viernes de 7:00 am a 3:00 pm.
  \item \textbf{Recepción y Conformidad:} Otorgada por el propietario en un plazo máximo de siete (7) días hábiles tras la culminación.
  \item \textbf{Garantía:} Mínimo 12 meses sobre la mano de obra y materiales. Responsabilidad por vicios ocultos de un (1) año.
  \item \textbf{Penalidades:} Penalidad por mora aplicable automáticamente por cada día de retraso injustificado.
  \item \textbf{Subcontratación y Adelantos:} No aplican.
\end{itemize}

% ===================================================================
\section{Especificaciones técnicas mínimas}
% ===================================================================

\subsection{Sistema de Puesta a Tierra}

\begin{tabular}{L{4.5cm}L{10cm}}
\toprule
\textbf{Componente} & \textbf{Especificación Técnica} \\
\midrule
Resistencia Exigida & No mayor a 5 ohmios ($\Omega$). Medición con telurómetro calibrado. \\
Electrodo & Varilla de cobre electrolítico, mínimo 5/8" de diámetro y 2.40 m de longitud (tipo Copperweld). \\
Conductor de Bajada & Cobre desnudo, calibre mínimo \#4 AWG. \\
Tratamiento de Suelo & Sustrato químico reductor de resistividad, no corrosivo y de larga duración (ej. Thor Gel). \\
Material de Relleno & Sal industrial de alta pureza o tierra de cultivo seleccionada. \\
Conexiones & Abrazaderas de cobre o soldadura exotérmica (cadweld). \\
Caja de Inspección & De concreto o polipropileno de alta resistencia con tapa registrable. \\
\bottomrule
\end{tabular}

\subsection{Sistema de Bombeo (Electrobomba)}

\begin{tabular}{L{4.5cm}L{10cm}}
\toprule
\textbf{Componente} & \textbf{Especificación Técnica} \\
\midrule
Potencia & \textbf{2 HP}. \\
Tipo & Electrobomba centrífuga para agua limpia. \\
Accesorios & Suministro e instalación de válvulas de compuerta, válvula check, niples, conexiones y tuberías de succión/impulsión en PVC pesado. \\
Tablero de Control & Suministro e instalación de tablero eléctrico con protección (llave termomagnética, contactor, relé térmico/guardamotor) y flotadores de nivel. \\
\bottomrule
\end{tabular}

\subsection{Grupo Electrógeno de Respaldo}

\begin{tabular}{L{4.5cm}L{10cm}}
\toprule
\textbf{Componente} & \textbf{Especificación Técnica} \\
\midrule
Capacidad & Dimensionado para una Máxima Demanda de \textbf{6.10 kW / 10 kVA} (régimen Prime/Standby). \\
Suministro & Grupo electrógeno cabinado e insonorizado, Tablero de Transferencia Automática (TTA) y sistema de escape. \\
Puesta en servicio & Incluye primera dotación de combustible, aceite, refrigerante y pruebas de funcionamiento en vacío y con carga. \\
\bottomrule
\end{tabular}

% ===================================================================
\section{Experiencia del proveedor y personal}
% ===================================================================

\begin{itemize}
  \item \textbf{Experiencia del Proveedor:} Acreditar un mínimo de dos (02) servicios similares de instalación electromecánica o puesta a tierra en los últimos tres (03) años.
  \item \textbf{Técnico Responsable:} Experiencia mínima de tres (03) años en instalaciones eléctricas y montaje.
  \item \textbf{Profesional Certificador:} Certificado de operatividad emitido y firmado por un Ingeniero Electricista o Mecánico Electricista colegiado y habilitado.
\end{itemize}

\vspace{1cm}

\begin{center}
\rule{8cm}{0.4pt} \\
Renzo Gabriel Mamani Galindo \\
Propietario
\end{center}

\end{document}
